\subsection{Path A: Stratum-Month Discrete-Time Hazard Model}\label{sec:patha}

Path A is the aggregate corroboration leg: a stratum-month discrete-time hazard fit to Freddie Mac loan-level performance data. It recovers \$928.9 billion in aggregate, 121.5\% standalone---reported as corroboration only and excluded from every headline figure. Appendix~\ref{app:patha-detail} collects its estimation detail, diagnostics, and cross-correlation exhibits.

The same caution applies to the paths themselves: Path B's $+0.190$ lag-0 levels correlation is carried by the window's shared trend ($-0.23$ differenced, $-0.20$ linearly detrended). The ABM gives $-0.20$ and $-0.30$ (from $-0.32$ in levels), Path A $-0.21$ and $-0.29$ (from $-0.44$): detrending collapses them into a cluster between $-0.20$ and $-0.30$, none positive, so Path B's distinction is its level accuracy, not its timing. Against the i.i.d.\ $1.96/\sqrt{n}$ reference ($\pm 0.30$ at $n = 42$ detrended, $\pm 0.31$ at $n = 41$ differenced), Path A ($-0.29$) and Path B ($-0.20$) sit inside the band under both empirical-series conventions. The ABM's detrended value straddles its edge ($-0.301$ common-series, $-0.303$ own back-out). Indistinguishability verdicts therefore rest on Table~\ref{tab:estimators}'s wider moving-block bands, under which none of the three is distinguishable from zero. (In levels, trend-contaminated and reported only for completeness, the simulated path co-moves most strongly with the rate from four months earlier, the empirical path contemporaneously.) No single channel carries the offset: the $\beta_B = 0$ and $\beta_1 = 0$ ablations both leave the peak at $-3$; shifting the rate input one to three months either way never reaches lag 0 ($-1$ at best, under a two-to-three-month lead), moving recovered trapped liquidity under \$1.5 billion across the entire scan. The offset is therefore smooth fitted dynamics against a seasonal empirical series (mechanism not isolated); Path A's seasonal re-estimation moving its peak from $-3$ to $-2$ (Section~\ref{sec:patha}) is consistent with seasonality carrying part of it. Figure~\ref{fig:ccf} plots the full cross-correlations. At the cohort level the residual is structural, not compositional. Across the six 30-year cohorts spanning the SOMA low-coupon mass (2.0--4.5\%), every simulated CPR tracks the contemporaneous rate change near-mechanically: first-difference correlation between $-0.87$ and $-0.94$, five of six phase-locked at lag~0, the 4.5\% coupon peaking at lag $-6$ with a weaker $r=-0.29$, reproducing the aggregate $-0.318$. By contrast, no realized cohort speed responds distinguishably from zero at any lag within $\pm 6$ months. Those realized speeds are measured independently in the Freddie 2017--2021 performance panel and the Fannie replication cohorts, with Freddie lag-0 first-difference correlations $-0.05$ to $-0.16$, Fannie $-0.03$ to $-0.14$, all inside the $\pm 2/\sqrt{n}$ band, maximum magnitude $0.34$ over all cohorts and lags, and no coherent sign or lag. The aggregate empirical non-response is therefore not a composition artifact, nor the simulated lead a datable cohort-level phase shift. One honest scope limit travels with this exhibit: realized cohort-level monthly speeds exist only in the loan-level agency panels anchoring Path B, not in the SOMA holdings the empirical CPR is backed out of, so the cohort split runs in the Path B universe, and the estimator-level timing falsification above stays at the SOMA aggregate. A cumulative counterpart completes it. Monthly differences being seasonality-dominated, the cross-section is read in cumulated levels. Age-matched Freddie realized prepayment runs from 9.43\% annualized CPR (28.1\% of balance) in the shallowest surviving gap bucket ($[-1,0)$ points, mean coupon 6.02\%) to 5.24\% (16.4\%) at $-3$ points and deeper (mean coupon 3.37\%). That is a gradient of $+4.20$ CPR points, four-way stratum-cluster 95\% interval $[+3.59, +4.66]$, one-sided permutation $p$ 0.005. This is the one realized-data exhibit in this paper that moves in the lock-in direction. It is not corroboration of the imported elasticity's magnitude: the production hazard implies $+0.94$ points at the central $\beta_1$ across the same buckets, so the realized gradient is roughly $4.5\times$ the model's. The excess carries everything travelling with coupon: credit composition, within-window refinancing on high-coupon shallow buckets, seasoning past the age cut. Detection power is 0.68: the design could have resolved $\beta_1$'s implication, and what it found was larger and confounded. The composition limb is tested: standardizing the deep bucket to the shallow's vintage $\times$ FICO $\times$ LTV composition (22 cells, 18 of them carrying both endpoint buckets) leaves the gradient at $+3.44$ points (95\% CI $[+1.17, +5.14]$, imputed weight share $0.3\%$), an exact decomposition putting $+0.77$ of the raw $+4.20$ on composition and $+3.43$ within common cells (run \texttt{episode\_confrontation\_within}). Against the implied $+0.94$, the excess is not composition in dimensions the panel can hold fixed, loan age now among them: adding 12-month age bands to the standardization cell (29 cells, the same 18 in common; run \texttt{episode\_gradient\_age\_bands}) leaves the gradient at $+3.37$ points, and the same cut on age alone puts $+0.38$ of the raw $+4.20$ --- $9.1\%$ --- on seasoning composition, age in the cell moving total composition from $+0.77$ to $+0.85$. The age cut itself is held at production, so this stratifies above the cut rather than testing where it sits. The interval moves against the exhibit even as the point survives: on the age-augmented cell it runs $[+0.60, +4.35]$ (imputed weight share $0.5\%$), covering the implied $+0.94$ where the committed $[+1.17, +5.14]$ excluded it. On run \texttt{episode\_confrontation\_within}'s own pre-committed branch map, that moves this configuration from a gap that persists to one that closes. So the exhibit weakens without collapsing: seasoning composition is measured and small, and the gradient is no longer separated from the model's implication by its own sampling interval. I had pre-committed that at least $80\%$ of the primary selection's shallow exposure would sit in single-band vintages; the realized share is $60.6\%$: age is not merely a relabelling of vintage in this book, so the control is stronger than I expected. That $+0.94$ sits at the in-sample floor and production ramp, so I recomputed it across the headline calibration (run \texttt{episode\_gradient\_recompute}): the shortfall is not an artifact of the demoted anchor. At the $4.991\%$ floor the implied gradient is $+0.76$ under the production form and $+0.89$ under the additive one, so the realized-to-implied ratio is $5.5$ and $4.7$ against the $4.5$ at the in-sample anchor, and across the floor-by-form-by-ramp grid every finite ratio lies between $3.6$ and $12.7$. One cell is not finite: at the headline floor on a $75$~PSA ramp both endpoint buckets are fully censored, so the model implies \emph{no} cross-sectional gradient and the comparison degenerates. I had pre-committed two signs and one was wrong: the max-form gradient fell as predicted, but the additive form lowered it too rather than raising it. A floor in survival space compresses proportional differences even where it never censors. I therefore read the cross-section as a signed, significant level gradient whose composition component is measured and small---consistent with lock-in operating, still unusable as a measurement of its size: within-window refinancing on high-coupon shallow buckets stays unresolved, no field here separating it from moving, and a level gradient is not a dollar marginal. The monthly-timing nulls above stand unchanged.